\section{Introduction}
\label{sec:intro}

\epigraph{Beauty is the first test: there is no permanent place in the world
for ugly mathematics.}{---G.~H.~Hardy, \emph{A Mathematician's Apology} (1940)}

The Efficient Market Hypothesis (EMH) remains the dominant theory of
price formation. It predicts when prices reflect available information
but does not provide a generative microstructural account of when and
why they fail to: bubbles, crashes, persistent cross-sectional return
anomalies, and the stylised observation that real markets very rarely
sit at the textbook supply--demand fixed point. Decades of empirical
work have produced rival theories whose claims often contradict one
another, because abundant and noisy financial data can be mined to
support nearly any hypothesis \citep{moffitt2017}. To avoid
contributing yet another empirical controversy, this paper takes the
less travelled \emph{atheoretical} route: constructing an economy from
the ground up in a controlled environment and asking what behaviour
emerges from a small number of common-sense rules.

The central technical claim of the paper is that two structural
features --- max-bid clearing and per-epoch feedback of clearing
prices into agent valuations --- generate persistent positive drift
in realised market prices even when individual agents' estimation
errors are exactly zero-mean. The mechanism is order-statistic bias:
each seller selects the maximum of multiple noisy bids, and the
maximum of $n$ zero-mean draws is, in expectation, strictly above
their mean. Feeding the realised maximum back into the next round's
anchor compounds the bias multiplicatively. No behavioural assumption
about investor optimism, risk aversion, or attention is required.
Bubbles, business cycles, persistent overshoots, and persistent
undershoots all follow from this single mechanism interacting with
heterogeneous time-on-market dynamics.

We use agent-based modelling \citep{wilensky1999} because it allows
this kind of structural mechanism to emerge from a few simple,
transparent rules. The setting is a real-estate market in a small
neighbourhood --- a market in which every participant must form a
personal estimate of an asset whose true value is fundamentally
uncertain. Real estate is a useful laboratory: prices are observable,
transactions are costly enough that na\"ive arbitrage arguments fail,
and divergence of opinion among buyers and sellers is the rule
rather than the exception.

The central object of study is the \emph{Estimated Dynamic Equilibrium
Model} (EDEM), in which supply and demand are not static schedules
but realisations of a stochastic process driven by heterogeneous,
error-prone agent valuations. EDEM is a dynamic extension of
\citet{miller1977}'s divergence-of-opinion theory, but it relaxes
Miller's static one-shot setting to allow valuations, the price
level, and the population of buyers and sellers to co-evolve over
time.

The paper makes three contributions:
\begin{enumerate}
\item \textbf{An order-statistic mechanism for bubbles without
biased agents} (\cref{sec:framework-asymmetry}): we derive in closed
form, for a clean special case, that the expected winning bid
exceeds the home value by $\sigma(n-1)/(n+1)$ when $n$ bidders
draw zero-mean errors with dispersion $\sigma$. Compounded across
epochs, this produces the exponential price drift of
\cref{sec:exp-bubble}.

\item \textbf{A unified framework} (\cref{sec:framework}) that nests
the classical Walrasian equilibrium and Miller's static premium as
limiting cases, and admits a family of out-of-equilibrium regimes
(business cycles, persistent overshoots and undershoots, bubbles,
constant transition) parameterised by estimation noise, balancer
strength, and time-on-market.

\item \textbf{Eight controlled experiments}
(\cref{sec:experiments}) implemented in the Python ABM framework
Mesa \citep{masad2015mesa}, replacing earlier
\citet{wilensky1999} NetLogo prototypes. The Mesa codebase is open
source and reproduces every figure in the paper from a clean
checkout.
\end{enumerate}

A practical implication is that empirical premium-vs-discount
findings depend strongly on the sample's position along the
underlying stochastic trajectory, which short-window studies cannot
identify. Two studies that find opposite signs may both be correct
characterisations of distinct epochs in the same process. A
secondary implication concerns valuation algorithms: a
machine-learning estimator trained on historical clearing prices
inherits the bid-selection asymmetry and re-deploys it as a
positively-biased point estimator, providing the market with a
coordinating signal that can amplify the very bubble the algorithm
was meant to price. \Cref{sec:disc-ml} discusses this in connection
with the 2021 wind-down of Zillow Offers \citep{zillow2021ibuying}.

The remainder of the paper is organised as follows. \Cref{sec:background}
reviews \citet{miller1977} and the divergence-of-opinion debate.
\Cref{sec:framework} presents the EDEM framework formally and derives
the order-statistic drift. \Cref{sec:implementation} describes the
Mesa implementation. \Cref{sec:experiments} reports the eight
experiments and their results. \Cref{sec:discussion} discusses
implications. \Cref{sec:extensions} sketches extensions and
\cref{sec:conclusion} concludes.
